\section{Consequences for a fixed number of layers}\label{sec:corolarios}

We obtained the limit of $f_h(n)$ by comparing two consecutive layers. To pass from layer $h$ to the full board, it suffices to
multiply the ratios between successive layers. Since there are only finitely many layers in the product, we can also sum their
logarithmic errors.

\begin{proof}[Proof of Corollary~\ref{cor:cociente}] By definition~\eqref{eq:fh}, $N_j/N_{j-1}=(n-j+1)/(j^2f_j)$ for each $j\ge1$.
Multiplying over $j=1,\dots,h$ gives \[
 \frac{N_h}{Q}=\prod_{j=1}^h\frac{n-j+1}{j^2f_j}=\frac{(n)_h}{(h!)^2\prod_{j=1}^hf_j},
\] and Theorem~\ref{thm:limite} gives $\prod_{j\le h}f_j\to F(0)^h>0$ and $(n)_h\sim n^h$. We have simply iterated the exact definition
of $f_h$. \end{proof}

\begin{proof}[Proof of Corollary~\ref{cor:capas}] For each $h\le K$, $f_h(n)\to F(0)$ by Theorem~\ref{thm:limite}, and
$F((h-\frac12)/n)\to F(0)$ and $F(h/n)\to F(0)$ by continuity of $F$ at zero (Theorem~\ref{thm:perfil}). Since $F(0)>0$, each of the
$K$ summands tends to zero. \end{proof}

The proof sums only $K$ terms and does not justify interchanging the limits $K\to\infty$ and $n\to\infty$.
